% Discussion
% Generated by report/md2tex.py from Report.md -- edit Report.md, then re-run.

\section{Discussion}

\subsection{What the three tiers say together}

The four experiments are deliberately different in shape, and the comparison is the report's main technical claim: "parallelise it" is not one decision but three, and which one applies is determined by the workload's dependency structure, not by the hardware. A scan has no dependencies and scales until the serial merge bites. An iterative solver has one dependency per iteration and stops scaling when the iteration becomes shorter than the barrier. An irregular mining phase has data-dependent cost and stops scaling when the slowest shard determines the finish, which no amount of workers fixes. The same machine produced 6.93$\times$, 1.67$\times$ and 4.34$\times$ on three workloads of comparable size.

\subsection{An implementation defect worth reporting}

The first version of the scan tier created a fresh process pool for each of the four stages. Measured that way, the pipeline peaked at 2.62$\times$ at 4 workers and \emph{slowed down} at 16: the two smallest files spent more time spawning sixteen interpreters than parsing their contents, and the row-count-identical work took 63 s at 16 workers against 42 s at 4. Reusing one pool across all four stages, and spilling the read-only reference set to disk so each worker loads it once instead of receiving it with every chunk task, raised the same measurement to 6.93$\times$ at 16 workers. The defect was invisible in the algorithmic analysis and only appeared because the benchmark reported per-stage timings rather than one total: the large stage was scaling fine while the small ones were going backwards.

\subsection{Bottlenecks}

Three bottlenecks were predicted in the proposal and all three were confirmed. The serial gunzip (33.8 s) caps the end-to-end benefit of every later stage; the parent-side merge caps the scan tier; and the inter-iteration barrier caps PageRank. A fourth was not predicted: process-based parallelism multiplies memory, which is the binding constraint on the hardware this project targets. A fifth is a property of the machine rather than of the algorithm: the 16-worker configuration exceeds the 14 physical cores, and eight of those cores are efficiency cores with roughly half the throughput of the performance cores. Work split evenly by count is therefore not split evenly by time, and the slowest worker sets the finish. The 43\% efficiency at 16 workers consequently mixes two effects \textemdash{} the serial merge and heterogeneous core placement \textemdash{} and this experiment does not attempt to separate them.

\subsection{Threats to validity}

\textbf{The working set is a release market, not an industry.} As Section 4.1 documents, no field in this release identifies a production country, and the three definitions we tested select 25,997, 15,715 or 0 titles. The analysis is therefore about films released in Hong Kong, and statements about "the local industry" refer to the Cantonese-language bloc identified in Section 4.3, not to a curated list of Hong Kong productions. Any claim about market share or output volume would inherit this error; claims about network structure within a bloc are less exposed, because bloc membership is itself estimated from the collaboration data.

\textbf{Billing truncation.} Taking the eight top-billed acting credits per title systematically drops unbilled supporting players and background performers, whose absence inflates the apparent density of the core. We report results at top-3 and top-8 as a sensitivity check; the bloc structure is stable between them, but node counts differ by a factor of 2.4 (21,912 versus 51,878), so no absolute density or degree figure should be read as a property of the industry rather than of the truncation rule.

\textbf{IMDb coverage.} Cast lists are incomplete for pre-1970 films and for Cantonese productions whose credits were never transcribed. The pre-1980 peak in Figure 3 is partly a cataloguing artefact, and the 1980s trough is almost certainly one. Temporal conclusions are therefore stated as within-sample trends, not as industry history.

\textbf{Sampling in betweenness.} With 64 pivots the ranking of the top 50 nodes is preserved against a 256-pivot run, but individual scores carry sampling error and are reported only as ranks, never as magnitudes.
